(1) 令 $\mathbf n=(\sin\theta,0,\cos\theta)$，Pauli 恒等式 $(\mathbf n\cdot\boldsymbol\sigma)^2=I$。因此指数可精确求和为 $U(t)=\cos(\Omega t/2)I-i\sin(\Omega t/2)(\sin\theta\,\sigma_x+\cos\theta\,\sigma_z)$，不需要微扰或绝热近似。
(2) 由 $\langle\downarrow_z|\sigma_x|\uparrow_z\rangle=1$、$\langle\downarrow_z|\sigma_z|\uparrow_z\rangle=0$，翻转振幅为 $-i\sin\theta\sin(\Omega t/2)$。故 $P_\downarrow(t)=\sin^2\theta\,\sin^2(\Omega t/2)$，其值介于零和一之间且随倾斜角改变振荡深度。
(3) 两级测量概率相加为一，故 $\langle\sigma_z\rangle=1-2P_\downarrow=\cos^2\theta+\sin^2\theta\cos(\Omega t)$。演化中 $H$ 不含显式时间，因而 $d\langle H\rangle/dt=0$；初态的 $\langle H\rangle=(\hbar\Omega/2)\cos\theta$，全程保持该值，虽然 $\langle\sigma_z\rangle$ 可随时间振荡。
(4) $\theta=0$ 或 $\pi$ 时初态恰为新哈密顿量本征态，$P_\downarrow=0$；$\theta=\pi/2$ 时 $P_\downarrow=\sin^2(\Omega t/2)$，可以完全翻转。$\Omega t=2\pi$ 时 $U=-I$，所有测量概率与初态相同，但自旋态获得负的全局相位，这与旋转二分之一自旋的双值表示一致。
